\documentclass[12pt, a4paper]{article}

\usepackage[utf8]{inputenc}
\usepackage{parskip}
\usepackage[margin=2.5cm]{geometry}
\usepackage[danish]{babel}
\usepackage{lmodern}
\usepackage{booktabs}
\usepackage{float}
\usepackage{pdfpages}
\usepackage[T1]{fontenc}
\usepackage{hyperref}
\usepackage{todonotes}
\usepackage[danish=guillemets]{csquotes}

\title{
  Projektbeskrivelse\\
  {\large eksamensprojekt programmering b}
}

\author{Jeppe Bøgeskov og Alexander Schade}

\begin{document}

\maketitle
\section*{Valg af programmeringssprog}
Til projektet anvendes Rust.

\section*{Link til repo}
\url{https://github.com/ZBC-Slagelse-HTX-X/DDU-Pro-eksamensprojekt}

\section*{Problemformulering}
\paragraph{Obs.} Produktet, som programmeringsprojektet afføder, skal agere et subkomponent (\textbf{hvor selve programmeringsrelevante aspekter og overvejelser, så er segregeret fra den DDU-relevante aflevering}) i et større spil, der behandles i faget DDU (Teknikfag A).

Her indgår både et shopping-system til progression og spilverden består primært af et åbent hav, hvor forskellige fjender interagere med hinanden og spilleren (altså spillets entiteter).
\begin{displayquote}
  \begin{enumerate}
    \item Hvordan kan man lave et system, der approksimerer, hvordan havstrømme påvirker andre entiteter på en optimal måde---her forstås, at beregninger er forholdsvis billige.
    \item Desuden undersøges det, hvordan man kan dataficere varer og kode logikken til et butikssystem.
    \item Afrundingsvis hvordan opstilles en pæn og klar grænseflade mellem brugger grænsefladen og data-logikken?
\end{enumerate}
\end{displayquote}

\section*{Ansvarsfordeling}
Det forventes, at Alexander Schade har ansvaret for forståelse af spilmotoren og integrering. Jeppe Bøgeskov holder styr på den overordnet udvikling af systemet.

\end{document}
